\section{Introduction}
\label{sec:intro}

The exponent \(\omega\) of matrix multiplication is the infimum of the
real numbers \(a\) such that two \(n\times n\) matrices can be multiplied
with \(O(n^a)\) arithmetic operations. In a preprint dated 2~October 2026
and released on 6~October, \citet{OAI26} proved \(\omega\le9/4\) over the
complex numbers, well below the bounds near \(2.371\) obtained by the
laser method~\citep{AD25,Dup26,Zad26,OAI26Stag}. The proof ends with an
existence theorem for tensor characters, and specifies neither a size
\(n\) nor an algorithm for \(n\times n\) matrices with fewer than
\(n^{9/4+\eps}\) multiplications.

The preceding paper~\citep{FinCon26} constructs, for every rational
\(\eps>0\), explicit exact bilinear algorithms with rational coefficients
and exponent below \(9/4+\eps\). Their sizes are doubly exponential in
\(1/\eps\), necessarily so within that construction: as \(\eps\to0\),
the least source side \(A\) (defined below) of a scheme with exponent at most
\(9/4+\eps\) satisfies \(\log\log A=\Theta(1/\eps)\), and its seed sides
are doubly exponential too~\citep[Corollary~11.2]{FinCon26}. The second
open question of~\citep[Section~11.5]{FinCon26} asks for a finite
construction with exponent below \(9/4+\eps\) in which \(\log A\) is
bounded by a polynomial in \(1/\eps\). That question was posed for the
source side \(A\), but \(A\) alone is not a robust measure: composition
with the classical algorithm reduces it to~\(1\)
(Remark~\ref{rem:composition}). The robust measures are the copy counts
\(r_*\) and \(m_*\) and the seed side \(n_0\)
(Section~\ref{subsec:intro-results}), which are also doubly exponential
within~\citep{FinCon26} (Table~\ref{tab:intro-comparison}). Here all
three are singly exponential. Our schemes have
\(\log r_*,\log m_*=O(\eps^{-9/2}\log^3(1/\eps))\), the least possible
order within the construction (Theorems~\ref{thm:intro-main}
and~\ref{thm:intro-limits}), and seeds with
\(\log n_0=O(\eps^{-11/2}\log^3(1/\eps))\)
(Theorem~\ref{thm:intro-seed}). They also have
\(\log A=O(\eps^{-9/2}\log^3(1/\eps))\), which answers the question as
posed.
Sections~\ref{sec:stages}--\ref{sec:limits} prove
Theorems~\ref{thm:intro-main}--\ref{thm:intro-limits} below, with long
proofs in Appendices~\ref{sec:lower-proof} and~\ref{sec:coefficients}.

\subsection{Results}
\label{subsec:intro-results}

We use the conventions of~\citep{FinCon26}, recalled in
Section~\ref{sec:prelim} below; all logarithms are natural. We write
\(T\restr T'\) if linear substitutions on each leg of the tensor \(T\) transform it into \(T'\),
\(\MM_d\) for \(d\times d\) matrix multiplication, and \(mT\) for the
direct sum of \(m\) copies of \(T\). A \emph{relative
scheme}~\citep[Definition~2.3]{FinCon26} is a restriction
\(r_*\MM_A\restr m_*\MM_B\) with \(B>A\); it computes \(m_*\) products
of \(B\times B\) matrices from \(r_*\) products of \(A\times A\)
matrices. We call \(A\) and \(B\) its \emph{source} and \emph{target
sides}, \(r_*\) and \(m_*\) its \emph{copy counts}, and
\(p=\log(r_*/m_*)/\log(B/A)\) its \emph{pooled exponent}. Applied
recursively, with the subproblems of each depth pooled into batches of
\(m_*\), it multiplies \(n\times n\) matrices with \(O(n^p)\) arithmetic
operations if \(p>2\), over every field that contains its
coefficients~\citep[Lemma~2.4]{FinCon26}.

Our schemes use the convolution tensor
\(C(r,b)=\sum_{\alpha<r}\sum_{\beta<b}x_\alpha y_\beta z_{\alpha+\beta}\),
which multiplies polynomials with \(r\) and \(b\) coefficients, the dot
product \(B_X(d)=x\sum_{\gamma<d}y_\gamma z_\gamma\), and the product
\(Q_{r,b}\) of the six copies of \(C(r,b)\) with permuted legs; note that
\(Q_{1,t}\cong\MM_{t^2}\). A \emph{stage} \((r,b)\to(c,r)\) of arity
\(d\) turns a power of \(Q_{r,b}\) and auxiliary dot products
\(B_X(d^N)\) into the same power of \(Q_{c,r}\) and a matrix
multiplication tensor. A \emph{chain} applies stages in sequence, from a
\emph{source pair} \((r_1,b_1)\) to a \emph{terminal pair} \((1,t)\)
(Section~\ref{subsec:intro-idea}).
\begin{theorem}[Main theorem]\label{thm:intro-main}
For every rational \(0<\eps\le1\), an explicit chain of at most
\(240\,\eps^{-3/2}\) stages, whose parameters are computed from \(\eps\)
by integer arithmetic, gives a relative scheme
\(r_*\MM_A\restr m_*\MM_B\) with rational coefficients such that
\(9/4<p<9/4+\eps/2\),
\[
 \log m_*<1.8\cdot10^{10}\,\eps^{-9/2}\log^3\frac{4096}\eps,\qquad
 \log r_*<2.7\cdot10^{10}\,\eps^{-9/2}\log^3\frac{4096}\eps,
\]
\(m_*=B^3\), \(r_*<A^{35/4}\), \(B\le A^2\) and
\(\log A<3\cdot10^9\,\eps^{-9/2}\log^3(4096/\eps)\). For
\(0<\eps\le\frac12\), the scheme with \(2\eps\) in place of \(\eps\)
satisfies \(9/4<p<9/4+\eps\),
\(\log m_*<8.4\cdot10^8\,\eps^{-9/2}\log^3(2048/\eps)\),
\(\log r_*<1.23\cdot10^9\,\eps^{-9/2}\log^3(2048/\eps)\) and
\(\log A<1.4\cdot10^8\,\eps^{-9/2}\log^3(2048/\eps)\).
\end{theorem}

\begin{theorem}[Seeds and one algorithm]\label{thm:intro-seed}
Let \(0<\eps\le1\) be rational, and let \(A\) be the source side of the
scheme of Theorem~\ref{thm:intro-main}.
\begin{enumerate}
\item[(a)] An explicit integer \(n_0\) with
\(\log n_0<2.1\cdot10^{11}\,\eps^{-11/2}\log^3(4096/\eps)\) has the
following property: for every \(n\ge n_0\), unrolling the recursion of
the scheme at side \(n\) gives a rational bilinear algorithm for
\(n\times n\) matrices with fewer than \(n^{9/4+\eps}\)
multiplications.
\item[(b)] The coefficients of the scheme are fractions
\(\mathfrak c/\mathcal D\) with a common denominator \(\mathcal D\),
where \(\mathcal D\) and all \(|\mathfrak c|\) are below
\(2^{\beta_\eps}\) with
\(\log_2\beta_\eps<10^4\,\eps^{-3/2}\log^2(4096/\eps)\). If
\(\eps=u/v\) with positive integers \(u,v\le A\), which holds whenever
\(\eps=1/v\) (then \(v<A\)), a generator computes \(\mathcal D\) and all numerators
\(\mathfrak c\) from \(u\) and \(v\) with at most \(A^{46}\) bit
operations.
\item[(c)] One algorithm multiplies \(n\times n\) matrices over every
field of characteristic zero, for every \(n\ge2\), in
\(n^{9/4+\delta(n)}\) steps, where
\(\delta(n)=O((\log\log n)^{6/11}(\log n)^{-2/11})\). A step is a step of
a fixed multitape Turing machine during the preprocessing, or a field
operation.
\end{enumerate}
\end{theorem}

Theorem~\ref{thm:intro-main} restates Theorem~\ref{thm:main} and
Corollary~\ref{cor:main-eps}, with Lemma~\ref{lem:resources}(a) for the
number of stages; the bounds on \(\log m_*\) and \(\log r_*\) follow
from \(m_*\le A^6\) and \(r_*<A^{35/4}\), as in
Remark~\ref{rem:limits-orders}. Parts (a), (b) and (c) of
Theorem~\ref{thm:intro-seed} are Theorem~\ref{thm:seed}, where
\(\log n_0\ge(24/\eps)\log A\), Proposition~\ref{prop:coefficients} and
Theorem~\ref{thm:uniform}. All these sizes are astronomically large, and
no algorithm here is shown to use \(O(n^{9/4})\) operations. The local
maps of Section~\ref{sec:stages} hold over every field; the relations
that use the separator and exactification are proved over \(\Q\), and we
do not treat positive characteristic.

For a fixed target, much smaller chains suffice.
Proposition~\ref{prop:witnesses}, restated next, certifies some with
the supplementary script \path|shifted_chain_certificates.py|, through
the \emph{fixed-batch bound}
\(\Phi=\log(2r_*/m_*)/\log(B/A)=p+\log2/\log(B/A)\)
of~\citep[(10.1)]{FinCon26}, which also bounds the exponent of the
recursion without pooling; Theorem~\ref{thm:intro-main} holds with
\(\Phi\) in place of \(p\) (Theorem~\ref{thm:main}).

\begin{theorem}[Finite examples; computer-assisted]
\label{thm:intro-examples}
Explicit chains give relative schemes over \(\Q\) with \(\Phi<2.5\),
\(2.3\), \(2.26\), \(2.84063\) and \(2.67283\), whose source sides have
\(33{,}379\), \(12{,}371{,}380\), \(48{,}859{,}824{,}830\), \(280\) and
\(4{,}653\) decimal digits, respectively. The last two targets lie about
\(10^{-5}\) below the floating-point values of \(\Phi\) for the schemes
of~\citep[Table~2]{FinCon26} with \(a=8\) and \(a=12\)
(Section~\ref{subsec:examples-comparison}).
\end{theorem}

Section~\ref{sec:limits} asks what the construction cannot do. Its
\emph{chain schemes} (Definition~\ref{def:class}) allow crops and
exchanges of the pairs, several blocks with their own widths and powers
in each stage, and every \emph{supply} of the auxiliary matrix
multiplication tensors, by the source or by tensors that earlier stages
produce (\emph{forward recycling}, Lemma~\ref{lem:recycling}).

\begin{theorem}[Limits within the construction]\label{thm:intro-limits}
\leavevmode
\begin{enumerate}
\item[(a)] Every chain scheme has \(p>9/4\), also with compressed
auxiliaries of any rank (Definition~\ref{def:compressed}).
\item[(b)] Let \(0<\eps<1/9216\). Every chain scheme with
\(p\le9/4+\eps\) satisfies
\(\log r_*>\log m_*>3.9\cdot10^{-6}\,\eps^{-9/2}\log^3(1/\eps)\) and
\(\log(B/A)>9.8\cdot10^{-5}\,\eps^{-4}\log^2(1/\eps)\). If the source
supplies every auxiliary, then also
\(\log A>1.3\cdot10^{-6}\,\eps^{-9/2}\log^3(1/\eps)\).
\end{enumerate}
\end{theorem}

Part (a) is Corollary~\ref{cor:compression-limit}, which extends
Lemma~\ref{lem:profile-potential}(c) by a Hankel flattening. Part (b) is
Theorem~\ref{thm:lower}. Its bounds on \(r_*\), \(m_*\) and
\(\log(B/A)\) extend to composites and tensor products of chain schemes
and classical algorithms, and Theorem~\ref{thm:intro-main} attains their
orders (Remark~\ref{rem:limits-orders}). The source side depends on the
supply. Forward recycling keeps the \(r_*\), \(m_*\) and
\(p<9/4+\eps/2\) of that scheme, changes \(\log(B/A)\) by at most
\(\log2\), and gives \(\log A<5.32\cdot10^6\,\eps^{-3}\log^3(4096/\eps)\)
(Corollary~\ref{cor:recycled-main}). It lowers \(\log n_0\) of the
seed only by a factor below \(1.34\), since the copy counts, not the
source side, set the size of the seed (Remark~\ref{rem:limits-orders}).
With computer assistance, Proposition~\ref{prop:small-sources} shows that every chain scheme with
\(p<3\), also after composition and tensor products, has
\(m_*>2^{494}\), and that the least such \(m_*\) is attained. None of
these limits is a lower bound for \(\omega\) or for tensor rank.

\begin{table}[ht]
\centering
\small
\begin{tabular}{@{}llll@{}}
\toprule
 & \multicolumn{2}{l}{Preceding paper} & This paper\\
\cmidrule(r){2-3}
 & size & source & size\\
\midrule
least \(\log m_*\) and \(\log r_*\) & \(>(6/5)^{9/(2\eps)}\)
 & \citep[Prop.~11.1(b)]{FinCon26}
 & \(\Theta(\eps^{-9/2}\log^3(1/\eps))\)\\
\(\log n_0\) of a seed & \(>\frac12(6/5)^{9/\eps}\)
 & \citep[Cor.~11.2(c)]{FinCon26}
 & \(O(\eps^{-11/2}\log^3(1/\eps))\)\\
least \(\log A\) & \(e^{\Theta(1/\eps)}\)
 & \citep[Cor.~11.2(b)]{FinCon26}
 & \(\Theta(\eps^{-9/2}\log^3(1/\eps))\)\\
\(\delta(n)\) of one algorithm & \(O(1/\log\log n)\)
 & \citep[Rem.~9.8]{FinCon26}, sketch
 & \(O((\log\log n)^{6/11}(\log n)^{-2/11})\)\\
digits of \(A\), \(\Phi\approx2.6728\) & \(234{,}905{,}651{,}862\)
 & \citep[Table~2]{FinCon26}, \(a=12\) & \(4{,}653\)\\
digits of \(A\), \(p<2.26\) & \(>9.29\cdot10^{34}\)
 & Section~\ref{subsec:examples-comparison} & \(48{,}859{,}824{,}830\)\\
\bottomrule
\end{tabular}
\caption{The two constructions, each with the source supplying every
auxiliary; far from \(9/4\), other constructions are smaller
(Section~\ref{subsec:examples-comparison}). The first three rows concern
exponents at most \(9/4+\eps\) as \(\eps\to0\): schemes with
\(p\le9/4+\eps\), and seeds with fewer than \(n_0^{9/4+\eps}\) products.
Preceding paper: the first and third rows hold for \(\eps\le1/4\), and
the second for \(\eps\le1/2\). In the first row,
\citep[Proposition~11.1(b)]{FinCon26} bounds \(\log A\) and gives
\(a\ge4\) for the target index \(a\); with this,
\citep[(7.3), (7.6) and~(11.2)]{FinCon26} and \(B>A\) give
\(\log r_*>\log m_*>2\log A\). This paper:
Theorems~\ref{thm:intro-main} and~\ref{thm:intro-limits}(b) for the
first and third rows (Remark~\ref{rem:limits-orders}),
Theorem~\ref{thm:intro-seed} for the second and fourth, and
Theorem~\ref{thm:intro-examples} for the last two.}
\label{tab:intro-comparison}
\end{table}

\subsection{The idea}
\label{subsec:intro-idea}

\paragraph{One stage.}
Fix \(d\ge2\) and \(1\le c<b\), and put \(h=b-c\). The \emph{shifted
sum} \((g_0,\ldots,g_{d-1})\mapsto\sum_{j<d}\xi^{jh}g_j\) maps \(d\)
polynomials with \(b\) coefficients onto those with \(b^+=db-(d-1)c\)
coefficients, with kernel the direct sum of \(d-1\) copies of the
polynomials with \(c\) coefficients. As it commutes with multiplication
by a polynomial, \(C(r,b)B_X(d)\) is block triangular in unimodular
bases adapted to the kernel, with a \emph{quotient} part \(C(r,b^+)\)
and \(d-1\) \emph{kernel} parts \(C(r,c)\) sharing the first leg
(Lemma~\ref{lem:shifted-chain}). If \((d+2)c\le db-r+1\), a crop and a
reset cut the quotient part into two \emph{outer} parts \(C(r,c)\) and
one \emph{middle} part \(C(r,c)^{(23)}\) (Lemma~\ref{lem:reset}). In the
\(N\)-th tensor power, prescribing the numbers \((\nK,\nO,\nmid)\) of
kernel, outer and middle parts on legs \(2\) and \(3\) deletes every
cross term, since each goes from a part to one of larger weight (the
weights are \(2\), \(1\), \(0\), in this order), and leaves
\(M=N!\,(d-1)^{\nK}2^{\nO}/(\nK!\,\nO!\,\nmid!)\) sectors sharing the
first leg (Lemma~\ref{lem:stage-types}), with
\((d+2)^N/(N+1)^2\le M<(d+2)^N\) for the largest \(M\)
(Lemma~\ref{lem:types}). L\"u's version~\citep{Lu26} of the separator
of~\citet{OAI26} turns \(2M-1\) copies of this sum into \(M\) independent sectors, each with an extra
dot product \(B_X(M)\). In the six orientations of the legs, the dot
products \(B_X(e)\) form \(\MM_{e^2}\), and the middle parts only
permute the orientations, so Proposition~\ref{prop:stage} gives a
degeneration over \(\Q\), of degree at most \(6(M-1)\),
\[
 (2M-1)^6\,Q_{r,b}^{\otimes N}\,\MM_{d^{2N}}\degen
 M^6\,Q_{c,r}^{\otimes N}\,\MM_{M^2} ,
\]
where \(\MM_e\) denotes \(e\times e\) matrix multiplication and \(M\)
the sector count. For the largest \(M\), the stage trades an auxiliary
of side \(d^{2N}\) for a matrix multiplication tensor of side \(M^2\approx(d+2)^{2N}\), at
the price of fewer than \(64\) copies.

\paragraph{Chains, and why \(9/4\).}
Interpolation~\citep{Bini80}, in the rational form
of~\citep[Lemmas~2.2 and~9.1(c)]{FinCon26}, makes each stage exact at a
tensor power \(k\), with \(E=1+6k(M-1)\) times as many copies
(Corollary~\ref{cor:exact-stage}). The power of \(Q_{c,r}\) in the
target of a stage is the power of \(Q\) in the source of the next, so
exact stages compose. Unlike~\citep{FinCon26}, a chain needs no
balancing of frequencies and no \emph{catalyst}, that is, no factor
common to both sides of a relation that each use of the relation
consumes and reproduces~\citep[Proposition~7.1]{FinCon26}. A chain, started by evaluating
\(Q_{r_1,b_1}\) with \(s^6\) products, \(s=r_1+b_1-1\), gives \(r_*\MM_A\restr m_*\MM_B\) with
\(A=\prod_id_i^{2Nk}\) and \(B=t^{2Nk}\prod_iM_i^{2k}\)
(Proposition~\ref{prop:chain-counts}). As \(N\) and \(k\) grow, with
types of the largest \(M\), \(p\) approaches the \emph{ideal exponent} \(3\log s/(\log t+G)\), where
\(G=\sum_i\log\frac{d_i+2}{d_i}\) is the \emph{gain}. The least profile
\(F\) (Section~\ref{subsec:prelim-profiles}) satisfies
\(dF(r,b)\ge(d+2)F(c,r)\) at every stage, \(F(1,t)=t\) and
\(F(r_1,b_1)\le s^{4/3}\), so \(\log t+G\le\frac43\log s\) and \(p>9/4\)
(Lemma~\ref{lem:profile-potential}). Equality holds on the \emph{arity
ramp}, the stages of arities \(3j-1\) and \(3j\) from \((R,R)\) down to
\((1,1)\), with gain \(\log F(R,R)\ge\frac43\log R\)
(Lemma~\ref{lem:ramp}). Above it, a \emph{bulk} of stages of one arity
\(D\) multiplies the pair by about the positive root \(\mu\) of
\(D\mu^2-\mu-(D+2)=0\) at each stage, with ideal exponent at most
\(9/4+27/(32(D+1)^2)\) (Lemmas~\ref{lem:bulk-potential}
and~\ref{lem:ideal-gap}). Theorem~\ref{thm:intro-main} takes
\(D=\lceil2\eps^{-1/2}\rceil\), \(R=8D\), about \(64D/\eps\) bulk stages
and \(N=k=O(\eps^{-3/2}\log(1/\eps))\), so that the excess of the ideal
exponent over \(9/4\) and all losses add up to less than \(\eps/2\) (Step~4 of the proof of Theorem~\ref{thm:main}). As \(A\)
depends only on the arities, \(N\) and \(k\),
\(\log A\le2N^2J\log(24D)=O(\eps^{-9/2}\log^3(1/\eps))\), where \(J\)
is the number of stages. The exponentially large source pair enters
only \(r_*\), through \(s^{6Nk}\), and the copy counts satisfy
\(m_*\le A^6\) and \(r_*<A^{35/4}\).

\subsection{Related work and credits}
\label{subsec:intro-related}

The bound \(\omega\le9/4\), the profile axioms and the convolution
relations behind them are due to \citet{OAI26}. The bound was extended
to every field~\citep{Nav26a,Zui26} (in Lean), to associative unital
rings~\citep{Nav26b}, to algebraically closed fields of every
characteristic, with a polynomial separator and the other steps checked
in outline~\citep{Lu26}, and to rectangular products in characteristic
zero~\citep{Uzn26}. The criterion of~\citet{Fil26} turns finite
families of shared-leg block identities with a nonnegative certificate
into exponent bounds, and also gives \(\omega\le9/4\) over every field.
These works prove that algorithms exist; to our knowledge, apart from~\citep{FinCon26}, no work gives
explicit sizes for exponents near \(9/4\). An OpenAI
preprint~\citep{OAI26Comp} proves \(\omega<2.258\) in characteristic
zero by a different method.

For \(h=1\) the shifted sum is known: for \(d=2\) it is the map in the
proof of~\citep[Lemma~4.1]{OAI26}, and for every \(d\) its kernel is the
first step \(\mathcal F^1\) of the determinant
filtration~\citep[Lemma~2.5]{Maakestad23},~\citep[Section~4.1]{FinCon26}
(there for \(b\ge d\)).
For \(h\ge2\) it is a substitution of binary forms
(Section~\ref{subsec:shifted-chain}); this is elementary, but we found
no source for it. New, to our knowledge, are keeping the kernel whole,
as \(d-1\) equal parts, the shift \(h\ge2\), the acyclic schedule
without a catalyst, the arity ramp, and the explicit bounds proved
here. The reset is
the partition in the proof of~\citep[Lemma~4.2]{OAI26} (see
also~\citep[Lemma~2.1(3)]{Uzn26}), and our type restriction
follows~\citep[Lemma~4.3]{FinCon26}. Shared-leg separation is due
to~\citet[Proposition~3.1]{OAI26}, with \(5M\) source copies; we use the
polynomial separator of~\citet[Proposition~2.1]{Lu26}, with \(2M-1\)
copies and degree at most \(M-1\), in the rational form
of~\citep[Lemma~9.1(b)]{FinCon26}. The ramp attains the diagonal
recursion~\citep[(5.1)]{OAI26} with equality, and
\(F(R,R)\ge R^{4/3}\) is~\citep[Lemma~5.1]{OAI26}. With \(N=1\) and the
cross terms removed by a degeneration, a stage gives a shared-leg block
identity \(C(r,b)B_X(d)\degen\sum_iT_i\) with \(d+2\) blocks, in the
sense of~\citet[Theorem~1.1, identity~(1.2)]{Fil26}, whose criterion
gives exponents but no sizes. Relative schemes are a form of relative
bilinear complexity~\citep{Str87}, and type restrictions are standard
in the laser method~\citep{Str87,CW90}.
